%% LyX 2.4.4 created this file.  For more info, see https://www.lyx.org/.
%% Do not edit unless you really know what you are doing.
\documentclass[english]{article}
\usepackage[T1]{fontenc}
\usepackage[utf8]{inputenc}
\setcounter{secnumdepth}{-2}
\setlength{\parindent}{0bp}
\usepackage{amsmath}
\usepackage{amsthm}
\usepackage{amssymb}

\makeatletter
\@ifundefined{date}{}{\date{}}
%%%%%%%%%%%%%%%%%%%%%%%%%%%%%% User specified LaTeX commands.
\usepackage{graphicx}
\usepackage{geometry}
\newcommand{\limi}{\underset{n\rightarrow\infty}{\lim}}
\newcommand{\limxz}{\underset{x\rightarrow x_{0}}{\lim}}
\newcommand{\limfxz}{\underset{x\rightarrow x_{0}}{\lim}f(x)}
\newcommand{\limfxm}{\underset{x\rightarrow x_{0}^{-}}{\lim}f(x)}
\newcommand{\limfxp}{\underset{x\rightarrow x_{0}^{+}}{\lim}f(x)}
\newcommand{\limxm}{\underset{x\rightarrow x_{0}^{-}}{\lim}}
\newcommand{\limxp}{\underset{x\rightarrow x_{0}^{+}}{\lim}}
\newcommand{\sqrm}{M_{m\times m}(\mathbb{F})}
\newcommand{\seqa}{(a_{n})_{n=1}^{\infty}}
\newcommand{\seqx}{(x_{n})_{n=1}^{\infty}}
\newcommand{\funcdr}{f:D\rightarrow\mathbb{R}}
\newcommand{\der}{\limxz\frac{f(x)-f(x_{0})}{x-x_{0}}}
\usepackage[all]{pdfprivacy}

\makeatother

\usepackage{babel}
\begin{document}
\title{{\Large Semester 2026A, Exam A, Q2\vspace{-2em}}}

\maketitle
\textbf{Problem.} Let $D\colon M_{n\times n}(\mathbb{F})\to\mathbb{F}$
be a volume function, and let $A\in M_{n\times n}(\mathbb{F})$.

We will denote by $\varepsilon$ the elementary row operation $R_{i}\to R_{i}+cR_{j}$,
where $1\le i,j\le n$, $i\ne j$, and $c\in\mathbb{F}$.

Prove that $D(\varepsilon(A))=D(A)$.

\medskip{}

\rule[0.5ex]{1\columnwidth}{1pt}

\medskip{}

\begin{proof}
Let us mark $A$ with $\begin{pmatrix}R_{1}\\
\vdots\\
R_{j}\\
\vdots\\
R_{i}\\
\vdots\\
R_{n}
\end{pmatrix}$ (we assume without loss of generality that $j<i$).

Then the matrix $A$ after operation $\varepsilon$ will be 

\[
\begin{pmatrix}R_{1}\\
\vdots\\
R_{j}\\
\vdots\\
R_{i}+cR_{j}\\
\vdots\\
R_{n}
\end{pmatrix}=\begin{pmatrix}R_{1}\\
\vdots\\
R_{j}\\
\vdots\\
R_{i}\\
\vdots\\
R_{n}
\end{pmatrix}+\begin{pmatrix}R_{1}\\
\vdots\\
R_{j}\\
\vdots\\
cR_{j}\\
\vdots\\
R_{n}
\end{pmatrix}
\]

Since a volume function is in particular a multi-linear function,

\[
D\left(\begin{pmatrix}R_{1}\\
\vdots\\
R_{j}\\
\vdots\\
R_{i}+cR_{j}\\
\vdots\\
R_{n}
\end{pmatrix}\right)=D\left(\begin{pmatrix}R_{1}\\
\vdots\\
R_{j}\\
\vdots\\
R_{i}\\
\vdots\\
R_{n}
\end{pmatrix}\right)+D\left(\begin{pmatrix}R_{1}\\
\vdots\\
R_{j}\\
\vdots\\
cR_{j}\\
\vdots\\
R_{n}
\end{pmatrix}\right)=D\left(\begin{pmatrix}R_{1}\\
\vdots\\
R_{j}\\
\vdots\\
R_{i}\\
\vdots\\
R_{n}
\end{pmatrix}\right)+cD\left(\begin{pmatrix}R_{1}\\
\vdots\\
R_{j}\\
\vdots\\
R_{j}\\
\vdots\\
R_{n}
\end{pmatrix}\right)
\]

Recall that for a matrix $B$ and for the elementary row operation
$R_{w}\leftrightarrow R_{k}$ which we mark with $\varepsilon'$,
$D\left(\varepsilon'\left(B\right)\right)=-D\left(B\right)$.

\[
\therefore D\left(\begin{pmatrix}R_{1}\\
\vdots\\
R_{j}\\
\vdots\\
R_{i}\\
\vdots\\
R_{n}
\end{pmatrix}\right)+cD\left(\begin{pmatrix}R_{1}\\
\vdots\\
R_{j}\\
\vdots\\
R_{j}\\
\vdots\\
R_{n}
\end{pmatrix}\right)=D\left(\begin{pmatrix}R_{1}\\
\vdots\\
R_{j}\\
\vdots\\
R_{i}\\
\vdots\\
R_{n}
\end{pmatrix}\right)+c\cdot-D\left(\begin{pmatrix}R_{1}\\
\vdots\\
R_{j}\\
R_{j}\\
\vdots\\
\vdots\\
R_{n}
\end{pmatrix}\right)
\]

And since $D$ is a volume function, it's an alternating function
in particular, which means it returns zero on a matrix that has two
adjacent identical rows: 
\[
c\cdot-D\left(\begin{pmatrix}R_{1}\\
\vdots\\
R_{j}\\
R_{j}\\
\vdots\\
R_{n}
\end{pmatrix}\right)=c\cdot-0=0
\]

Hence 

\[
D\left(\begin{pmatrix}R_{1}\\
\vdots\\
R_{j}\\
\vdots\\
R_{i}+cR_{j}\\
\vdots\\
R_{n}
\end{pmatrix}\right)=D\left(\begin{pmatrix}R_{1}\\
\vdots\\
R_{j}\\
\vdots\\
R_{i}\\
\vdots\\
R_{n}
\end{pmatrix}\right)
\]

Which means $D(\varepsilon(A))=D(A)$ as required. 
\end{proof}
\begin{quote}
This is an exam problem I translated and solved. The original exam
was written by the Linear Algebra 1 course lecturers at HUJI.
\end{quote}

\end{document}
